\section{Related Work}
\label{sec:related}

\textbf{FL autoencoders for IoT anomaly detection.}
Rey~et~al.~\cite{rey2022federated}, Khraisat~et~al.~\cite{khraisat2025fediot},
and Olanrewaju~et~al.~\cite{olanrewajugeorge2025flae}
apply federated autoencoders to device-partitioned N-BaIoT and report aggregate detection
metrics.
Thresholds are applied implicitly or globally; none isolates threshold calibration scope
as a variable or measures per-client FPR dispersion.

\textbf{Non-IID aggregation.}
Nguyen~et~al.~\cite{nguyen2025fedmse} improve AUC under Dirichlet non-IID N-BaIoT
by weighting aggregation (FedMSE); the approach modifies the training/aggregation protocol,
not the post-training threshold policy, and per-client alarm-rate dispersion is not reported.
General non-IID benchmarks~\cite{zhao2018noniid,li2022federated} focus on loss
and accuracy rather than per-client alarm-boundary dispersion.

\textbf{Model-level FL personalization.}
Personalization in FL is studied mainly at the model level~\cite{fallah2020personalized,arivazhagan2019federated,liang2020think}:
methods such as Per-FedAvg, FedPer, and LG-FedAvg modify model weights or training
objectives to adapt to local distributions.
These approaches adapt the model parameters; DATP instead fixes the shared AE and
adapts only the post-training alarm boundary---the two adaptation layers are complementary.

\textbf{Threshold-level personalization.}
Per-device thresholds are a natural deployment option acknowledged in
Kitsune~\cite{mirsky2018kitsune} and in DIoT~\cite{nguyen2019diot}, but neither provides
a controlled comparison fixing all other variables.
The existence of local thresholds is not new; the missing element is
a controlled study that fixes encoder, training protocol, and seeds across
threshold-scope variants and then measures per-client FPR dispersion.

Table~\ref{tab:related_comparison} summarizes the gap.

\begin{table}[t]
  \caption{Positioning: adaptation layer and threshold treatment.}
  \label{tab:related_comparison}
  \centering
  \resizebox{\columnwidth}{!}{%
  \renewcommand{\arraystretch}{1.2}%
  \begin{tabular}{lll}
    \toprule
    Work group & Adaptation layer & Threshold treatment \\
    \midrule
    FL-AE IoT IDS~\cite{rey2022federated,khraisat2025fediot} & Model (shared AE) & Implicit/global \\
    FedMSE~\cite{nguyen2025fedmse} & Aggregation weights & Not varied \\
    Per-FedAvg, FedPer~\cite{fallah2020personalized,arivazhagan2019federated} & Model weights/objectives & Not addressed \\
    Kitsune, DIoT~\cite{mirsky2018kitsune,nguyen2019diot} & Local AE & Per-device (no controlled study) \\
    \textbf{DATP (this work)} & \textbf{Threshold only} & \textbf{Controlled comparison} \\
    \bottomrule
  \end{tabular}}
\end{table}
